\documentclass[11pt]{scrartcl}

% ---------- Packages ----------
\usepackage[margin=1in]{geometry}
\usepackage{mathpazo}
\usepackage{amsmath,amssymb,amsthm}
\usepackage{tikz}
\usetikzlibrary{arrows.meta,positioning}
\usepackage[most]{tcolorbox}
\usepackage{xcolor}
\usepackage{hyperref}

% ---------- Colors ----------
\definecolor{accentblue}{RGB}{44,95,138}
\definecolor{boxfill}{RGB}{240,245,250}
\definecolor{boxborder}{RGB}{180,200,220}
\definecolor{highlightred}{RGB}{200,50,50}

% ---------- Theorem environments ----------
\newtcolorbox{thmbox}[1]{
  colback=boxfill, colframe=boxborder,
  boxrule=0.6pt, arc=2pt, left=6pt, right=6pt, top=4pt, bottom=4pt,
  fonttitle=\bfseries\color{accentblue}, title=#1
}
\theoremstyle{plain}
\newtheorem{theorem}{Theorem}
\theoremstyle{remark}
\newtheorem*{remark}{Remark}

\hypersetup{colorlinks=true, linkcolor=accentblue, urlcolor=accentblue}

% ---------- Footer ----------
\usepackage{scrlayer-scrpage}
\pagestyle{scrheadings}
\clearpairofpagestyles
\cfoot{\footnotesize Page~\thepage\ --- For error reporting or contributing, visit \url{https://github.com/nasro-bay/Graph_theory_proofs}}

\title{\color{accentblue}Hamiltonian Graphs and Digraphs}
\author{}
\date{}

\begin{document}
\maketitle

% =========================================================
\begin{thmbox}{{Theorem 3 (Ore, 1960)}}
If $G$ is a simple graph with $n\ (\geq 3)$ vertices, and $\deg(v) + \deg(w) \geq n$ for every pair of non-adjacent vertices $v$ and $w$, then $G$ is Hamiltonian.
\end{thmbox}

\begin{proof}
We assume the statement is false and derive a contradiction. So let $G$ be a non-Hamiltonian graph on $n$ vertices satisfying the given condition on the vertex degrees. By adding extra edges if necessary, we may assume that $G$ is ``just barely'' non-Hamiltonian, in the sense that adding any further edge produces a Hamiltonian graph (note that adding extra edges cannot violate the degree condition, since degrees can only increase). It follows that $G$ contains a path
$$v_1 \to v_2 \to \cdots \to v_n$$
passing through every vertex. But since $G$ is non-Hamiltonian, the vertices $v_1$ and $v_n$ are not adjacent, and so $\deg(v_1) + \deg(v_n) \geq n$; by the pigeonhole principle, there must be some vertex $v_i$ adjacent to $v_1$ with the property that $v_{i-1}$ is adjacent to $v_n$. But this gives the required contradiction, since
$$v_1 \to v_2 \to \cdots \to v_{i-1} \to v_n \to v_{n-1} \to \cdots \to v_{i+1} \to v_i \to v_1$$
is then a Hamiltonian cycle.

\begin{center}
\begin{tikzpicture}[scale=0.9, every node/.style={circle,draw=accentblue,fill=accentblue!15,minimum size=6mm,inner sep=0pt}]
  \node (v1)   at (0,0)    {$v_1$};
  \node (v2)   at (1.3,0)  {$v_2$};
  \node[draw=none,fill=none] (d1) at (2.6,0) {$\cdots$};
  \node (vim1) at (3.9,0)  {$v_{i-1}$};
  \node (vi)   at (5.3,0)  {$v_i$};
  \node (vip1) at (6.7,0)  {$v_{i+1}$};
  \node[draw=none,fill=none] (d2) at (8.0,0) {$\cdots$};
  \node (vnm1) at (9.3,0)  {$v_{n-1}$};
  \node (vn)   at (10.7,0) {$v_n$};

  \draw (v1)--(v2);
  \draw (vim1)--(vi)--(vip1);
  \draw (vnm1)--(vn);

  \draw[highlightred, thick] (v1) to[bend left=30] (vi);
  \draw[highlightred, thick] (vim1) to[bend left=55] (vn);
\end{tikzpicture}
\end{center}
\end{proof}

% =========================================================
\begin{thmbox}{{Corollary 7.2 (Dirac, 1952)}}
If $G$ is a simple graph with $n\ (\geq 3)$ vertices, and $\deg(v) \geq \dfrac{n}{2}$ for each vertex $v$, then $G$ is Hamiltonian.
\end{thmbox}

\begin{proof}
The result follows from the preceding theorem, since for all $v, u \in V(G)$ we have $\deg(v) + \deg(u) \geq n$, and hence $G$ is Hamiltonian.
\end{proof}

% =========================================================
\begin{thmbox}{Theorem (Hamiltonian Digraphs --- Ghouila-Houri)}
Let $D$ be a strongly connected digraph with $n$ vertices. If $\operatorname{outdeg}(v) \geq \dfrac{n}{2}$ and $\operatorname{indeg}(v) \geq \dfrac{n}{2}$ for each vertex $v$, then $D$ is Hamiltonian.
\end{thmbox}

\end{document}
